Fix \(\mathbb P\in\mathcal D_{d,\epsilon}^{\mathrm{obs}}\).  Proposition \(\mathrm{prop:identification\mbox{-}and\mbox{-}extension}\) gives
\[
 0\leq\Psi(\mathbb P)=\sum_x\bar f_\epsilon(q_x)
 \leq\sum_xs(q_x)=1
\]
and supplies a completion \(P\in\mathcal P_{d,\epsilon}\) whose observed margin is \(\mathbb P\).  Every statistic below depends only on the observed sample and the displayed auxiliary variables, so its law under the completion is its law under \(\mathbb P\).  Hence \(\mathrm{lem:centered\mbox{-}factorial\mbox{-}pilot\mbox{-}control}\) applies uniformly over the full observed model.

First suppose \(d\geq D_0\) and \(n\geq d/L_d\), and put \(m=n/8\).  In the uncapped Poisson comparison experiment, counts belonging to distinct \(x\)'s are independent by Poisson splitting; therefore the clipped statistics \(T_x^{\mathrm{JF}}\), including their pilot components, are independent.  The cell bias bound, \(\sum_xp_x=1\), and Cauchy--Schwarz give
\[
 \left|\sum_xE_{\mathbb P}T_x^{\mathrm{JF}}-\Psi(\mathbb P)\right|^2
 \leq C_\epsilon\left\{
 \frac d{mL_d}+\frac{d^2}{m^2L_d^2}\right\}.                     \tag{15}
\]
Independence and the cell variance bound give
\[
 \operatorname{Var}_{\mathbb P}\left(\sum_xT_x^{\mathrm{JF}}\right)
 \leq C_\epsilon d^{1/16}
 \left\{\frac{L_d}{m}+\frac{dL_d^2}{m^2}\right\}.               \tag{16}
\]
Because \(m\geq d/(8L_d)\), the second term of (15), divided by \(d/(mL_d)\), is at most eight.  The two terms of (16), divided by \(d/(mL_d)\), are bounded respectively by
\[
 d^{-15/16}L_d^2,
 \qquad 8d^{-15/16}L_d^4,
\]
which are uniformly bounded for integer \(d\geq2\).  Thus
\[
 E_{\mathbb P}\left[
 \left\{\sum_xT_x^{\mathrm{JF}}-\Psi(\mathbb P)\right\}^2\right]
 \leq C_\epsilon\frac d{mL_d}.                                    \tag{17}
\]
Projection onto \([0,1]\) cannot increase squared distance from \(\Psi(\mathbb P)\in[0,1]\).

Adjoin an infinite independent continuation of the fixed sample.  The ideal uncapped Poisson statistic and the capped randomized statistic in \(\mathrm{def:jackson\mbox{-}factorial\mbox{-}estimator}\) agree on \(\{M\leq n\}\).  Both projected losses are at most one, whereas exponential Markov gives
\[
 \Pr\{M>n\}=\Pr\{\operatorname{Pois}(n/4)>n\}\leq e^{-cn}.       \tag{18}
\]
Equations (17)--(18) bound the randomized fixed-sample risk by
\[
 C_\epsilon\frac d{nL_d}+e^{-cn}.
\]
For \(d\geq2\), the exponential term is at most a universal multiple of \(d/(nL_d)\).  Conditional Jensen for squared loss then shows that the Rao--Blackwellized estimator
\[
 \widehat V_{n,d}^{\mathrm{JF}}=E(\widetilde V\mid O_1,\ldots,O_n)
\]
has no larger risk.

If \(d\geq D_0\) and \(n<d/L_d\), the estimator is \(1/2\), its risk is at most \(1/4\), and \(d/(nL_d)>1\), so the minimum in the theorem equals one.

It remains to treat \(2\leq d<D_0\).  Conditional on \(\overline N_x=k>0\), the arm count \(\overline N_{ax}\) is binomial with success probability \(\rho_{ax}\geq\epsilon\).  For \(H\sim\operatorname{Bin}(k,\rho)\),
\[
 E\frac{\mathbf1\{H>0\}}H\leq\frac2{k\rho},
 \qquad
 \Pr(H=0)=(1-\rho)^k\leq\frac1{k\rho}.                            \tag{19}
\]
The first inequality follows from \(H^{-1}\mathbf1\{H>0\}\leq2/(H+1)\) and
\[
 E(H+1)^{-1}=\frac{1-(1-\rho)^{k+1}}{(k+1)\rho};
\]
the second follows from \((1-\rho)^k\leq(1+k\rho)^{-1}\).  Conditional binomial variance and the \(0/0=0\) convention imply
\[
 E\{(\widehat\mu_{ax}-\mu_{ax})^2\mid\overline N_x=k\}
 \leq\frac3{k\epsilon}.                                            \tag{20}
\]
The maximum is one-Lipschitz in the sup norm.  Weighted Cauchy--Schwarz followed by (20) yields
\[
 E\left[
 \left\{\sum_x\widehat p_x
 (\max_a\widehat\mu_{ax}-\max_a\mu_{ax})\right\}^2\right]
 \leq\frac{6d}{n\epsilon}.                                         \tag{21}
\]
Moreover,
\[
 E\lVert\widehat p-p\rVert_1^2
 \leq d\sum_x\operatorname{Var}(\widehat p_x)\leq\frac dn.       \tag{22}
\]
Decomposing the empirical-ratio error into the regression term in (21) and the mass-estimation term in (22), and also using that both the estimator and target lie in \([0,1]\), gives risk at most \(C_\epsilon\min\{1,d/n\}\).  Since \(L_d\) is bounded above and below for \(2\leq d<D_0\), this is at most
\[
 C_\epsilon\min\{1,d/(nL_d)\}.
\]
All bounds are uniform over \(\mathcal D_{d,\epsilon}^{\mathrm{obs}}\), so they permit unequal unknown propensities, arbitrary masses, null cells, overlap boundaries, outcome boundaries, and exact ties.  This proves the theorem in every regime.